\subsection{有限维 g-模的特征标}

令 \((e^{\lambda})_{\lambda \in \mathfrak{h}^{*}}\) 为环 \(\mathbf{Z}[\mathfrak{h}^*]\) 的典则基。在空间 \(\mathbf{Z}^{\mathfrak{h}^*}\) 上赋予所有从 \(\mathfrak{h}^*\) 到 \(\mathbf{Z}\) 的映射所成空间以各因子离散拓扑的积拓扑。若 \(\varphi \in \mathbf{Z}^{\mathfrak{h}^*}\)，则族 \((\varphi (\nu)e^{\nu})_{\nu \in \mathfrak{h}^*}\) 可求和，并且

\[
\varphi = \sum_ {\nu \in \mathfrak {h} ^ {*}} \varphi (\nu) e ^ {\nu}.
\]

令 \(\mathbf{Z}\langle\mathrm{P}\rangle\) 为这样的 \(\varphi\in\mathbf{Z}^{\mathfrak{h}^{*}}\) 所成的集合：其支集包含于有限个形如 \(\nu-\mathrm{P}_{+}\) 的集合之并，其中 \(\nu\in\mathfrak{h}^{*}\)。则 \(\mathbf{Z}[\mathrm{P}]\subset\mathbf{Z}\langle\mathrm{P}\rangle\subset\mathbf{Z}^{\mathfrak{h}^{*}}\)。在 \(\mathbf{Z}\langle\mathrm{P}\rangle\) 上定义一个扩展 \(\mathbf{Z}[\mathrm{P}]\) 中环结构的环结构，即对 \(\varphi,\psi\in\mathbf{Z}\langle\mathrm{P}\rangle\) 和 \(\nu\in\mathfrak{h}^{*}\)，置

\[
(\varphi \psi) (\nu) = \sum_ {\mu \in \mathfrak {h} ^ {*}} \varphi (\mu) \psi (\nu - \mu)
\]

（族 \((\varphi(\mu)\psi(\nu - \mu))_{\mu \in \mathfrak{h}^*}\) 的支集是有限的，这是因为 \(\varphi\) 和 \(\psi\) 的支集满足上述条件）。若 \(\varphi = \sum_{\nu} x_{\nu} e^{\nu}\) 且 \(\psi = \sum_{\nu} y_{\nu} e^{\nu}\)，则 \(\varphi \psi = \sum_{\nu, \mu} x_{\nu} y_{\mu} e^{\nu + \mu}\)。

令 \(\nu \in \mathfrak{h}^*\)。\(\nu\) 关于正根的一个分拆是一个族 \((n_{\alpha})_{\alpha \in \mathrm{R}_{+}}\)，其中 \(n_{\alpha}\) 为满足 \(\geq 0\) 的整数 \(\nu = \sum_{\alpha \in \mathrm{R}_{+}} n_{\alpha} \alpha\)。我们用 \(\mathfrak{P}(\nu)\) 表示 \(\nu\) 关于正根的分拆数。有

\[
\mathfrak {P} (\nu) > 0 \Longleftrightarrow \nu \in \mathrm{Q} _ {+}.
\]

本段中，用 K 表示 \(Z\langle P\rangle\) 中的以下元素：

\[
\mathrm{K} = \sum_ {\gamma \in \mathrm{Q} _ {+}} \mathfrak {P} (\gamma) e ^ {- \gamma}.
\]

现在回忆（第 VI 章第3节第3号，命题 2）：

\[
d = \prod_ {\alpha \in \mathrm{R} _ {+}} (e ^ {\alpha / 2} - e ^ {- \alpha / 2}) = \sum_ {w \in \mathrm{W}} \varepsilon (w) e ^ {w \rho}
\]

是 \(\mathbf{Z}[\mathrm{P}]\) 中的反不变元素。

引理 1。在环 \(\mathbf{Z}\langle \mathrm{P}\rangle\) 中，有 \(\mathrm{K}\) 。\(\prod_{\alpha \in \mathrm{R}_{+}}(1 - e^{-\alpha}) = \mathrm{Ke}^{-\rho}d = 1\)。事实上，

\[
\mathrm{K} = \prod_ {\alpha \in \mathrm{R} _ {+}} \left(e ^ {0} + e ^ {- \alpha} + e ^ {- 2 \alpha} + \dots\right)
\]

因此

\[
K e ^ {- \rho} d = \prod_ {\alpha \in \mathrm{R} _ {+}} \left(1 + e ^ {- \alpha} + e ^ {- 2 \alpha} + \dots\right) \prod_ {\alpha \in \mathrm{R} _ {+}} \left(1 - e ^ {- \alpha}\right) = 1.
\]

引理 2。令 \(\lambda \in \mathfrak{h}^*\)。模 \(Z(\lambda)\)（§6，第3号）有一个属于 \(\mathbf{Z}\langle P\rangle\) 的特征标，并且有 \(d.\) ch \(Z(\lambda) = e^{\lambda + \rho}\)。

令 \(\alpha_{1},\ldots ,\alpha_{q}\) 为 \(\mathbb{R}_+\) 中互不相同的元素。\(X_{-\alpha_1}^{n_1}X_{-\alpha_2}^{n_2}\dots X_{-\alpha_q}^{n_q}\otimes 1\) 构成 \(\mathrm{Z}(\lambda)\) 的一个基（§6，命题 6 (iii)）。对 \(h\in \mathfrak{h}\)，有

\[
\begin{array}{l} h. (X _ {- \alpha_ {1}} ^ {n _ {1}} X _ {- \alpha_ {2}} ^ {n _ {2}} \ldots X _ {- \alpha_ {q}} ^ {n _ {q}} \otimes 1) \\ \qquad = [ h, X _ {- \alpha_ {1}} ^ {n _ {1}} \ldots X _ {- \alpha_ {q}} ^ {n _ {q}} ] \otimes 1 + (X _ {- \alpha_ {1}} ^ {n _ {1}} \ldots X _ {- \alpha_ {q}} ^ {n _ {q}}) \otimes h. 1 \\ \qquad = (\lambda - n _ {1} \alpha_ {1} - \dots - n _ {q} \alpha_ {q}) (h) (X _ {- \alpha_ {1}} ^ {n _ {1}} \ldots X _ {- \alpha_ {q}} ^ {n _ {q}} \otimes 1). \end{array}
\]

因此，\(\mathrm{Z}(\lambda)^{\lambda - \mu}\) 的维数为 \(\mathfrak{P}(\mu)\)。这证明了 \(\operatorname{ch} \mathrm{Z}(\lambda)\) 有定义，属于 \(\mathbf{Z}\langle \mathrm{P} \rangle\)，并且

\[
\operatorname{ch} Z (\lambda) = \sum_ {\mu} \mathfrak {P} (\mu) e ^ {\lambda - \mu} = \mathrm{Ke} ^ {\lambda}.
\]

现在只需应用引理 1。

引理 3。令 M 为一个 g-模，它有一个特征标 ch(M)，其支集包含于有限个 \(\mu-P_{+}\) 之并。令 U 为 g 的包络代数，Z 为 U 的中心，\(\lambda_{0}\in h^{*}\)，\(\chi_{\lambda_{0}}\) 为相应的从 Z 到 k 的同态（§8，定理 2 的推论 1）。假设对所有 \(\S8\)\(z\in Z\)，\(z_{M}\) 都是比值为 \(\chi_{\lambda_{0}}(z)\) 的伸缩。令 \(D_{M}\) 为所有 \(\lambda\in W(\lambda_{0}+\rho)-\rho\) 所成的集合，使得 \(\lambda+Q_{+}\) 与 Supp(ch M) 相交。则 ch(M) 是所有  的 ch \(Z(\lambda)\) 的 Z-线性组合。\(\lambda\in D_{M}\)

若 \(\mathrm{Supp(chM)}\) 为空，引理显然成立。假设 \(\mathrm{Supp(chM)} \neq \emptyset\)。令 \(\lambda\) 为此支集的极大元素，并置 \(\dim M^{\lambda} = m\)。存在一个从  到 M 的 \(\mathfrak{g}\)-同态，它将 \(\varphi\)\((Z(\lambda))^{m}\)\((Z(\lambda)^{\lambda})^{m}\) 双射地映到 \(M^{\lambda}\)（§6，第3号，命题 6 (i)）。因此，\(\S 6\)\(Z(\lambda)\) 的中心特征标是 \(\chi_{\lambda_0}\)，所以 \(\lambda \in W(\lambda_0 + \rho) - \rho\)（§8，第5号，定理 2 的推论 1）。这证明 \(\S 8\)\(D_M \neq \emptyset\)，并可对 \(CardD_M\) 作归纳。令 L、N 分别为 \(\varphi\) 的核与余核。于是有如下 \(\mathfrak{g}\)-同态正合序列：

\[
0 \to \mathrm{L} \to (\mathrm{Z} (\lambda)) ^ {m} \to \mathrm{M} \to \mathrm{N} \to 0
\]

因此

\[
\operatorname{ch} (\mathrm{M}) = - \operatorname{ch} (\mathrm{L}) + m \operatorname{ch} \mathrm{Z} (\lambda) + \operatorname{ch} (\mathrm{N})
\]

（§7，第7号，公式 (6)）。Supp(ch L) 和 Supp(ch N) 包含于有限个 \(\mu -\mathrm{P}_{+}\) 之并。对 \(z\in \mathbf{Z}\)，\(z_{\mathrm{L}}\) 和 \(z_{\mathrm{N}}\) 都是比值为 \(\chi_{\lambda_{0}}(z)\) 的伸缩。显然，\(D_{N} \subset D_{M}\)。另一方面，\((\lambda + Q_{+}) \cap \text{Supp(ch M)} = \{\lambda\}\)，而 \(\lambda \notin \text{Supp(ch N)}\)，所以 \(\lambda \notin D_{N}\)，从而

\[
\mathrm{Card} \mathrm{D} _ {\mathrm{N}} <   \mathrm{Card} \mathrm{D} _ {\mathrm{M}}.
\]

另一方面，L 是 \((Z(\lambda))^{m}\) 的子模；若 \(\lambda^{\prime} \in D_{L}\)，则 \(\lambda^{\prime} + Q_{+}\) 与 Supp(ch L) ⊂ Supp ch Z(λ) 相交，所以 \(\lambda \in \lambda^{\prime} + Q_{+}\)（§6，第1号，命题 1 (ii)）；由此 \(D_{L} \subset D_{M}\)。由于 \(L \cap (Z(\lambda)^{\lambda})^{m} = 0\)，有 \(\lambda \notin D_{L}\)，所以

\(\mathrm{CardD_L} <   \mathrm{CardD_M}\)

现在只需应用归纳假设。

定理 1（H. Weyl 特征标公式）。令 M 为有限维单 g-模，\(\lambda\) 为其最高权。则

\[
\left(\sum_ {w \in \mathrm{W}} \varepsilon (w) e ^ {w \rho}\right). \mathrm{chM} = \sum_ {w \in \mathrm{W}} \varepsilon (w) e ^ {w (\lambda + \rho)}.
\]

沿用引理 3 的记号，M 的中心特征标为 \(\chi_{\lambda}\)（§6，第4号，命题 7）。因此，由引理 2 和 3，d.ch M 是满足下式的 \(\S6\)\(e^{\mu+\rho}\) 的 Z-线性组合：

\[
\mu + \rho \in W (\lambda + \rho).
\]

另一方面，由 §7，第7号，引理 7，d.ch M 是反不变的，且其唯一的极大项为 \(e^{\lambda+\rho}\)，故定理得证。

例。取 \(\mathfrak{g} = \mathfrak{sl}(2,k)\)，\(\mathfrak{h} = kH\)。令 \(\alpha\) 为 \((\mathfrak{g},\mathfrak{h})\) 的满足 \(\alpha (H) = 2\) 的根。\(\mathfrak{g}\)-模 \(V(m)\) 的最高权为 \((m / 2)\alpha\)。因此

\[
\begin{array}{r l} & {\mathrm{ch} (\mathrm{V} (m)) = (e ^ {(m / 2) \alpha + \frac {1}{2} \alpha} - e ^ {- (m / 2) \alpha - \frac {1}{2} \alpha}) / (e ^ {\frac {1}{2} \alpha} - e ^ {- \frac {1}{2} \alpha})} \\ & {\qquad = e ^ {- (m / 2) \alpha}. (e ^ {(m + 1) \alpha} - 1) / (e ^ {\alpha} - 1)} \\ & {\qquad = e ^ {- (m / 2) \alpha} (e ^ {m \alpha} + e ^ {(m - 1) \alpha} + \dots + 1)} \\ & {\qquad = e ^ {(m / 2) \alpha} + e ^ {(m - 2) \alpha / 2} + \dots + e ^ {- (m / 2) \alpha}} \end{array}
\]

这也可由 §1，第2号，命题 2 立即得到。

